\begin{center}
\begin{tikzpicture}[
    node distance = 1.2cm and 4cm,
    point/.style = {circle, fill, inner sep=1.8pt},
    blue_arr/.style = {blue, -{Stealth[scale=0.8]}, thin},
    red_arr/.style = {red, -{Stealth}, shorten >=5pt, shorten <=5pt},
    font=\large
]

    % --- 1. PROBABILITÀ E NODI DI INGRESSO ---
    \node (p1) at (-4, 1) {$P(x_1) = P_1 = p$};
    \node (p2) at (-4, -1) {$P(x_2) = P_2 = 1 - p$};

    \node (xl1) at (-0.6, 1) {$x_1$};
    \node[point] (x1) at (0, 1) {};
    
    \node (xl2) at (-0.6, -1) {$x_2$};
    \node[point] (x2) at (0, -1) {};

    % --- 2. NODI DI USCITA ---
    \node[point] (y1) at (4, 1) {};
    \node (yl1) at (4.6, 1) {$y_1$};
    
    \node[point] (y2) at (4, -1) {};
    \node (yl2) at (4.6, -1) {$y_2$};

    % --- 3. COLLEGAMENTI (Linee nere del canale) ---
    \draw (x1.center) -- (y1.center);
    \draw (x2.center) -- (y2.center);
    \draw (x1.center) -- (y2.center);
    \draw (x2.center) -- (y1.center);

    % --- 4. FRECCE ROSSE (Input) ---
    \draw[red_arr] (p1) -- (xl1);
    \draw[red_arr] (p2) -- (xl2);

    % --- 5. ANNOTAZIONI BLU (Probabilità di transizione) ---
    % Sopra x1 -> y1
    \draw[blue_arr] (1.5, 1.4) -- node[above, font=\small] {$1-\alpha$} (2.5, 1.4);
    
    % Sotto x2 -> y2
    \draw[blue_arr] (1.5, -1.4) -- node[below, font=\small] {$1-\alpha$} (2.5, -1.4);
    
    % Diagonale x1 -> y2 (freccia alpha)
    \begin{scope}[shift={(2.8,-0.35)}, rotate=-26]
        \draw[blue_arr] (0,0) -- node[right, font=\small] {$\alpha$} (0.8,0);
    \end{scope}
    
    % Diagonale x2 -> y1 (freccia alpha)
    \begin{scope}[shift={(2.5,0.1)}, rotate=26]
        \draw[blue_arr] (0,0) -- node[right, font=\small, pos=0.8] {$\alpha$} (0.8,0);
    \end{scope}

\end{tikzpicture}
\end{center}